\section{Resultados}

Este capítulo apresentará as evidências produzidas pela aplicação do protocolo avaliativo. Os resultados da prova de conceito serão utilizados apenas como referência preliminar e não serão confundidos com os resultados da execução final da dissertação.

\subsection{Caracterização da bateria e do conjunto de referência}

% TODO: quantidade por categoria, critérios de balanceamento, respondibilidade,
% processo de curadoria e estatísticas do gold standard.

\subsection{Resultados da recuperação documental}

% TODO: comparar BM25, busca vetorial e busca híbrida por Recall@1, Recall@5,
% Recall@10, MRR e nDCG@k. Incluir Precision@k e Hit@k somente se forem
% compatíveis com a granularidade real dos julgamentos, além de intervalos,
% diversidade, estratificações e resultados por categoria.

\subsection{Resultados da geração de respostas}

% TODO: apresentar dimensões da rubrica sem misturar observação e interpretação.

\subsection{Avaliação humana e avaliação automatizada}

% TODO: perfil e número dos avaliadores, piloto, concordância, estabilidade do
% juiz automatizado e divergências relevantes.

\subsection{Análise de erros e casos críticos}

% TODO: documentar falhas de recuperação, atribuição, escopo, extrapolação,
% insuficiência documental e casos negativos ou ambíguos.
